\section{Introducci\'on y motivaci\'on}

La apnea obstructiva del sue\~no (AOS) y las alteraciones de la arquitectura del sue\~no representan dos problemas cl\'inicos estrechamente relacionados. En la pr\'actica, ambos se estudian con polisomnograf\'ia (PSG), una modalidad rica en se\~nales pero costosa, instrumentada y dif\'icil de escalar. El proyecto del curso propone un escenario m\'as austero y, por ello, m\'as desafiante: resolver clasificaci\'on de etapas de sue\~no y detecci\'on de apnea utilizando \textbf{solo EEG monocanal}. Esa restricci\'on obliga a dise\~nar mejor el preprocesamiento, las \emph{features}, la validaci\'on y el an\'alisis cr\'itico.

En este trabajo se desarroll\'o la l\'inea principal del laboratorio con modelos cl\'asicos tabulares. En lugar de usar una arquitectura profunda como cuerpo principal, se privilegi\'o una bater\'ia reproducible basada en \emph{SVM-RBF}, \emph{Random Forest} y \emph{XGBoost}, porque el enunciado exige comparar familias tradicionales bajo el mismo protocolo y responder preguntas sobre desbalance, interpretabilidad, robustez y generalizaci\'on. Esa decisi\'on permite adem\'as vincular directamente los resultados con la literatura cl\'asica de \emph{power spectral density}, complejidad, entrop\'ia y clasificadores basados en \'arboles o kernels \cite{gao2023psdrf,nakamura2017complexity,li2022cascadedsvm}.

El estudio se organiz\'o en dos ejes experimentales. El primer eje fue \textbf{staging} con \expname{sleep_edf_expanded_tuned} como baseline intra-dataset y con dos pruebas \emph{cross-dataset} hacia/desde ISRUC. El segundo eje fue \textbf{apnea + staging multitarget} con \expname{mitbih_apnea_stage_classic} como baseline intra-dataset y con dos evaluaciones externas entre MIT-BIH PSG y St. Vincent. Esta partici\'on permite cubrir los requisitos acad\'emicos m\'inimos del laboratorio sin bono: validaci\'on \emph{subject-wise}, generalizaci\'on \emph{cross-dataset}, tres familias de modelos y comparaci\'on con baselines publicados.

La hip\'otesis operativa del trabajo es doble. Primero, que un conjunto compacto de \feature{eeg\_*} extra\'idas de ventanas de 30 s puede producir baselines razonables para staging y apnea si el protocolo es estrictamente \emph{subject-wise}. Segundo, que el mayor cuello de botella no ser\'a la optimizaci\'on dentro de un dataset sino la \textbf{transferencia entre cohortes}, debido a diferencias de canal equivalente, mezcla de sujetos, calidad de se\~nal, prevalencia de eventos y sesgo hacia clases mayoritarias. El resto del informe se organiza para poner a prueba esa hip\'otesis de forma reproducible y trazable a artefactos reales del repositorio.
